\documentclass[12pt,a4paper]{article}
\usepackage[utf8]{inputenc}
\usepackage[portuguese]{babel}
\usepackage{amsmath, amsfonts, amssymb}
\usepackage{graphicx}
\usepackage{booktabs}
\usepackage{listings}
\usepackage{xcolor}
\usepackage{geometry}
\usepackage{hyperref}
\usepackage{float}

\geometry{margin=1in}

% Configuração para exibição de código fonte (Python/MATLAB)
\definecolor{codegreen}{rgb}{0,0.6,0}
\definecolor{codegray}{rgb}{0.5,0.5,0.5}
\definecolor{codepurple}{rgb}{0.58,0,0.82}
\definecolor{backcolour}{rgb}{0.96,0.96,0.94}

\title{
  \vspace{-1cm}
  {\large ECI / Politécnica — UFRJ} \\[0.4em]
  {\large COC473 — Álgebra Linear Computacional} \\[0.8em]
  \rule{\linewidth}{0.5pt} \\[0.6em]
  {\LARGE \textbf{Método Iterativo dos
Gradientes Conjugados}} \\[0.4em]
  {\large Trabalho 2} \\[0.2em]
  \rule{\linewidth}{0.5pt}
}
\author{
  Leonardo Peres Albertazzi Drummond \\
  Marco Antonio Tronco Felix \\
  Eduardo de Mendonça Freire \\
  João Pedro Peçanha Cotrim \\
  Carlos Augusto Bertão Rodrigues \\
  Andre Pacífico
}
\date{\today}

\begin{document}


\maketitle

\section{Introdução}
O Método dos Gradientes Conjugados (CG -- \textit{Conjugate Gradient}) é um dos algoritmos iterativos mais potentes e eficientes para a resolução de sistemas de equações lineares de grande porte na forma:
\begin{equation}
    A x = b
\end{equation}
onde $A \in \mathbb{R}^{n \times n}$ representa uma matriz simétrica e definida positiva (SDP), $b \in \mathbb{R}^n$ é o vetor de termos independentes, e $x \in \mathbb{R}^n$ é o vetor de incógnitas a ser computado.

Ao contrário dos métodos diretos tradicionais como a Eliminação de Gauss ou a Fatoração de Cholesky, cuja complexidade computacional é de ordem $\mathcal{O}(n^3)$, o método dos Gradientes Conjugados converge teoricamente no espaço de no máximo $n$ passos (em aritmética exata). O custo por iteração é dominado pelo produto matriz-vetor, gerando uma complexidade aproximada de $\mathcal{O}(n^2)$ para matrizes cheias ou $\mathcal{O}(m)$ para matrizes esparsas estruturadas com $m$ elementos não nulos, tornando-o ideal para problemas físicos modelados por equações diferenciais.

\section{Implementação}
O método funciona gerando uma base de vetores ortogonais em relação ao produto interno definido pela matriz $A$ (direções conjugadas), ou seja, direções que satisfazem $v_i^T A v_j = 0$ para todo $i \neq j$. 

Abaixo encontra-se a modelagem matemática e lógica do algoritmo iterativo com base na sequência exata de cálculo apresentada no pseudocódigo de referência:

\subsection{Algoritmo Matemático}
\begin{enumerate}
    \item \textbf{Inicialização:}
    Dado um chute inicial $x_0$ (geralmente adotado como o vetor de aproximações iniciais):
    \begin{align*}
        r_0 &= b - A x_0 \quad (\text{Resíduo Inicial}) \\
        v_0 &= r_0 \quad (\text{Direção de busca inicial}) \\
        i &= 0
    \end{align*}
    
    \item \textbf{Laço Iterativo:} Realizado de forma iterativa até que um critério de parada pré-estabelecido seja satisfeito (como a norma do resíduo cair abaixo de uma tolerância $\epsilon$):
    \begin{align*}
        \alpha_i &= \frac{v_i^T r_i}{v_i^T A v_i} \quad (\text{Comprimento do passo}) \\[1.5ex]
        x_{i+1} &= x_i + \alpha_i v_i \quad (\text{Solução aproximada}) \\[1.5ex]
        r_{i+1} &= r_i - \alpha_i A v_i \quad (\text{Novo resíduo}) \\[1.5ex]
        \beta_i &= -\frac{v_i^T A r_{i+1}}{v_i^T A v_i} \quad (\text{Fator de correção / Novo comprimento de passo \&}) \\[1.5ex]
        v_{i+1} &= r_{i+1} + \beta_i v_i \quad (\text{Nova direção de busca}) \\[1.5ex]
        i &= i + 1
    \end{align*}
\end{enumerate}


\section{Matriz Exemplo e Casos de Teste}
Para validar o comportamento numérico e a taxa de convergência do algoritmo dos Gradientes Conjugados (CG), utilizou-se o experimento proposto no livro-texto. O caso de estudo consiste em uma matriz esparsa $A$ de dimensões $500 \times 500$, cuja construção foi realizada através das seguintes etapas:
\begin{enumerate}
    \item Atribuição do valor $1.0$ a todas as posições da diagonal principal ($a_{ii} = 1.0$).
    \item Preenchimento das posições fora da diagonal com números aleatórios extraídos de uma distribuição uniforme no intervalo $[-1, 1]$, garantindo a simetria inicial da matriz de modo que $A = A^T$.
    \item Aplicação de um filtro de esparsidade comandado pelo parâmetro de limiar $\tau$: qualquer elemento fora da diagonal principal que satisfaça $|a_{ij}| > \tau$ é substituído pelo valor zero.
\end{enumerate}

Como descrito na literatura teórica, para valores de $\tau$ muito próximos de zero, o sistema gera uma matriz bem-condicionada e estritamente positiva definida. À medida que o valor de $\tau$ aumenta, tanto a densidade de elementos não nulos quanto o número de condicionamento pioram, fazendo com que a matriz perca a propriedade de ser definida positiva para limiares mais altos (como $\tau = 0.2$).

\subsection{Análise das Propriedades da Matriz}
Antes da execução dos métodos numéricos, realizou-se uma análise estrutural e espectral da matriz gerada para cada valor do parâmetro $\tau$. A Tabela 1 documenta o número de condicionamento ($\kappa(A)$), a densidade de elementos não-nulos e a verificação matemática da condição de matriz Simétrica Positiva Definida (SPD).

\begin{table}[H]
\centering
\caption{Propriedades estruturais e espectrais da matriz $500 \times 500$ em função de $\tau$.}
\label{tab:propriedades_matriz}
\vspace{2mm}
\begin{tabular}{ccccc}
\toprule
\textbf{Limiar ($\tau$)} & \textbf{Condicionamento $\kappa(A)$} & \textbf{Simétrica?} & \textbf{PD?}  & \textbf{Diagonal Dominante?} \\
\midrule
$0.01$ & $1.0815$ & Sim & Sim & Sim\\
$0.05$ & $2.3932$ & Sim & Sim & Não \\
$0.10$ & $459.37$ & Sim & Não & Não \\
$0.20$ & $824.99$  & Sim & Não & Não \\
$0.50$ & $824.99$  & Sim & Não & Não \\
$1.00$ & $5560.5$  & Sim & Não & Não \\
\bottomrule
\end{tabular}
\end{table}

Os resultados empíricos obtidos demonstram que o aumento do limiar $\tau$ aumentam o número de condicionamento da matriz. Na qual, esse fator irá gerar impactos no fator de convergência dos métodos iterativos.

Além disso, a partir de $\tau = 0.10$ a matriz deixa de ser positiva definida e a partir de $\tau = 0.05$ ela deixa de ser diagonal dominante.

\subsection{Resultados e Comparação entre Métodos}
O sistema linear $Ax = b$ foi resolvido adotando-se um vetor $x$ aleatório, usando o valor de $Ax$ como valor de $b$ e tomando um vetor de aproximação inicial $x_0 = \mathbf{1}$. Foram confrontados os desempenhos do Método dos Gradientes Conjugados (CG), Método de Jacobi e Método de Gauss-Seidel sob uma tolerância estipulada de $\epsilon = 10^{-8}$ ou um limite máximo de $2000$ iterações. 

A Tabela 2 detalha o número de iterações necessárias para a convergência de cada método em seus respectivos cenários.

\begin{table}[h!]
\centering
\caption{Comparação do número de iterações para convergência ($\epsilon = 10^{-8}$).}
\label{tab:comparacao_metodos}
\vspace{2mm}
\begin{tabular}{cccc}
\toprule
\textbf{Cenário de Teste} & \textbf{Gradientes Conjugados} & \textbf{Método de Jacobi} & \textbf{Gauss-Seidel} \\
\midrule
Matriz $\tau = 0.01$ & 5 iterações & 7 iterações & 5 iterações\\
Matriz $\tau = 0.05$ & 12 iterações & 20 iterações & 12 iterações \\
Matriz $\tau = 0.10$ & 156 iterações & Divergiu & Divergiu \\
Matriz $\tau = 0.20$ & Parou no máximo de iterações & Divergiu & Divergiu \\
Matriz $\tau = 0.50$ & Parou no máximo de iterações & Divergiu & Divergiu \\
Matriz $\tau = 1.00$ & 1231 iterações & Divergiu & Divergiu \\
\bottomrule
\end{tabular}
\end{table}

Denota-se que a cada vez que aumentamos o limiar $\tau$ o número de iterações aumenta para convergência dos métodos iterativos. Com isso, causando a divergência dos métodos de Jacobi e Gauss-Seidel. 

No entanto, a matriz apresenta um comportamento interessante no valor de $\tau = 1$, o método dos Gradientes Conjugados volta a convergir dentro do limite máximo de iterações.

A Tabela 3 detalha o valor de erro relativo final na convergência ou divergência de cada método em seus respectivos cenários.

\begin{table}[H]
\centering
\caption{Comparação do erro relativo resultante ($\frac{\|b - Ax\|_2}{\|b\|_2}$.).}
\label{tab:comparacao_metodos_erro}
\vspace{2mm}
\begin{tabular}{cccc}
\toprule
\textbf{Cenário de Teste} & \textbf{Gradientes Conjugados} & \textbf{Método de Jacobi} & \textbf{Gauss-Seidel} \\
\midrule
Matriz $\tau = 0.01$ & $2.705\times 10^{-09}$ & $1.959\times 10^{-11}$  & $6.499\times 10^{-11}$\\
Matriz $\tau = 0.05$ & $9.154\times 10^{-09}$ & $2.360\times 10^{-09}$& $3.098\times 10^{-10}$ \\
Matriz $\tau = 0.10$ & $8.702\times 10^{-09}$  & $2.859\times 10^{94}$  & $3.960\times 10^{143}$\\
Matriz $\tau = 0.20$ & $2.212\times 10^{-03}$& INF & INF \\
Matriz $\tau = 0.50$ & $2.212\times 10^{-03}$ & INF & INF \\
Matriz $\tau = 1.00$ & $9.547\times 10^{-09}$ & INF & INF \\
\bottomrule
\end{tabular}
\end{table}

Comparando o cenário que todos os métodos convergiram, $\tau = 0.01$ e $\tau = 0.05$ vemos que os métodos de Jacobi e Gauss-Seidel apresentaram menor erro relativo que o Gradientes Conjugados, no entanto, analisando todos os resultados o Gradientes Conjugados mostrou-se mais eficiente, devido a apresentar convergência em 4 casos de 6 testes feitos.

\subsection{Gráficos de Log(r) x iterações}

\subsubsection{Matriz $\tau = 0.01$}
\includegraphics[width=0.9\linewidth]{gradientes_conjugados_0.01.png}
\subsubsection{Matriz $\tau = 0.05$}
\includegraphics[width=0.9\linewidth]{gradientes_conjugados_0.05.png}

Até $\tau  = 0.05$ a matriz é positiva definida, o que trouxe um comportamento de decrescimento linear para o gráfico.

\subsubsection{Matriz $\tau = 0.10$}
\includegraphics[width=0.9\linewidth]{gradientes_conjugados_0.1.png}

Diferente dos valores de $\tau$ anteriores esse gráfico retornou um comportamento não monótono, apresentando pequenas "barrigas", no entanto, acelera na descida final resultando na convergência do método.

Além disso, diferente das matrizes anteriores, a matriz com valor de $\tau = 0.1$ se mostrou não ser positiva definida o que afetou a maneira da convergência do método.
\subsubsection{Matriz $\tau = 0.20$}
\includegraphics[width=0.9\linewidth]{gradientes_conjugados_0.2.png}
\subsubsection{Matriz $\tau = 0.50$}
\includegraphics[width=0.9\linewidth]{gradientes_conjugados_0.5.png}

Os gráficos acima ilustram perfeitamente o comportamento de estagnação e oscilação caótica do Método dos Gradientes Conjugados quando aplicado à matriz indefinida com limiar $\tau = 0.50$, atingindo o limite configurado de 2000 iterações sem alcançar a convergência. Ao contrário do cenário estável visto anteriormente, a perda da propriedade definida positiva da matriz introduz autovalores negativos ou nulos, o que sabota a minimização consistente da energia do sistema. Como consequência, o algoritmo perde a sua trajetória ótima de descida e entra em um regime de busca desordenada, fazendo com que o resíduo fique "preso" e flutue de forma ruidosa em uma faixa alta de erro, majoritariamente entre $10^{-1}$ e $10^{-3}$.

\subsubsection{Matriz $\tau = 1.00$}
\includegraphics[width=0.9\linewidth]{gradientes_conjugados_1.png}

Diferente dos gráfico anteriores apesar de possuir um comportamento oscilatório caótico, o gráfico apresentou um comportamento decrescente, o que resultou na sua futura convergência.

\section{Comparação com as matrizes do trabalho 1}

Para comparar a performance do método em detrimento dos outros métodos iterativos realizamos o teste nas 12 matrizes utilizadas no trabalho anterior. Foram confrontados os desempenhos do Método dos Gradientes Conjugados (CG) sob uma tolerância estipulada de $\epsilon = 10^{-8}$ ou um limite máximo de $3000$ iterações. 

\begin{table}[H]
\centering
\caption{Comparação do número de iterações e do erro relativo nas matrizes do trabalho 1.}
\label{tab:comparacao_matrizes}
\vspace{2mm}
\begin{tabular}{ccccc}
\toprule
\textbf{Matriz} & \textbf{n}  & \textbf{erro relativo} & \textbf{Gradientes Conjugados} & \textbf{SPD?} \\
\midrule
\texttt{bcsstk01}& 48  & $3.894 \times 10^{-9}$ & 142 iterações & sim \\
\texttt{bfw62b}& 62  & $5.637 \times 10^{-9}$ & 30 iterações & não\\
\texttt{bcsstm02}& 66  & $2.501\times 10^{-16}$ & 12 iterações & sim \\
\texttt{bcsstk03}& 112  & $9.148 \times 10^{-9}$  & 444 iterações & sim\\
\texttt{bcsstk04}& 132  & $9.392 \times 10^{-9}$ & 512 iterações & sim \\
\texttt{bcsstm22}& 138  & $4.280 \times 10^{-9}$ & 55 iterações & sim \\
\texttt{bfw398b}& 398  & $9.140 \times 10^{-9}$  & 38 iterações & não \\
\texttt{494\_bus}& 494  & $9.681 \times 10^{-9}$ & 1124 iterações & sim \\
\texttt{dwb512}& 512  &  $3.732 \times 10^{-9}$  & 16 iterações & não \\
\texttt{nos6}& 675  & $8.519 \times 10^{-9}$ & 672 iterações & sim\\
\texttt{gr\_30\_30}& 900  & $7.547 \times 10^{-9}$ & 63 iterações & sim \\
\texttt{1138\_bus}& 1138 & $9.916 \times 10^{-9}$ & 2017 iterações & sim \\
\bottomrule
\end{tabular}
\end{table}

Denota-se que diferente do trabalho 1, o método iterativo dos Gradientes Conjugados convergiu para todas as matrizes apresentando um erro relativo médio na ordem de $10^{-9}$. Demonstrando assim a maior eficiência desse método em relação aos outros. 

Vale ressaltar que para as matrizes bcsstm02 e bcsstm22 que são diagonais ela não convergiu em duas iterações, diferente do que aconteceu com os métodos de Gauss-Seidel e Jacobi.

\section{Conclusão}
A investigação prática realizada neste relatório permitiu mapear com clareza o comportamento assintótico e os limites de estabilidade do Método dos Gradientes Conjugados (CG) em comparação aos métodos estacionários de Jacobi e Gauss-Seidel. Os experimentos empíricos evidenciaram que o parâmetro de limiar $\tau$ atua diretamente sobre as propriedades espectrais da matriz de teste, ditando o seu número de condicionamento e a preservação de características algébricas fundamentais, tais como a dominância diagonal e a positividade definida. 

Ficou demonstrada a sensibilidade dos métodos de Jacobi e Gauss-Seidel às restrições do critério de convergência por diagonal dominante. À medida que o aumento de $\tau$ povoou a matriz com coeficientes aleatórios fora da diagonal principal, ambos os métodos sofreram divergências catastróficas abruptas já a partir de $\tau \ge 0.10$, atingindo valores de erro relativo em escalas de \textit{overflow} numérico (INF). Em contrapartida, a formulação expandida do método de Gradientes Conjugados implementada neste trabalho revelou-se muito mais resiliente às perturbações espectrais devido ao cálculo dinâmico do fator de correção $\beta_i$, que força a ortogonalidade entre as direções de busca a cada iteração. Embora o método tenha estagnado sob oscilações caóticas nos cenários onde a matriz tornou-se estritamente indefinida ($\tau = 0.20$ e $\tau = 0.50$), ele evitou o colapso numérico imediato por explosão de resíduo observável nas técnicas estacionárias.

O fenômeno mais marcante do estudo manifestou-se no cenário limite de $\tau = 1.00$. Embora a matriz tenha atingido o seu maior número de condicionamento registrado ($\kappa(A) = 5560.5$) e se estruturado de forma densa e indefinida, o método dos Gradientes Conjugados foi capaz de reestabelecer a sua trajetória de descida e alcançar a convergência em $1231$ iterações sob um erro relativo na ordem de $10^{-9}$. Esse comportamento valida empiricamente propriedades associadas a grandes sistemas aleatórios descritos na literatura matemática, nos quais o acúmulo de componentes simétricas densas propicia um rearranjo dos autovalores residuais que, uma vez purgados nas fases iniciais de estagnação ruidosa, engatam um regime de superconvergência tardia na reta final do algoritmo.

Por fim, a validação do algoritmo sobre o banco de dados do Trabalho 1 consagrou o Método dos Gradientes Conjugados como uma ferramenta amplamente superior e robusta para problemas de engenharia e modelagem física. O algoritmo alcançou a convergência em $100\%$ das doze matrizes testadas --- incluindo malhas severamente esparsas e de grande porte como a \texttt{1138\_bus} --- mantendo erros relativos rigidamente controlados na ordem de $10^{-9}$ a $10^{-16}$. Esclareceu-se também que, para matrizes diagonais, o método requer um número de passos condicionado à quantidade de autovalores distintos no espectro da matriz, desmistificando a falsa superioridade das duas iterações observadas nos métodos estacionários, as quais decorriam de uma interrupção precoce por divergência e não de convergência real. O método dos Gradientes Conjugados consolida-se, portanto, como a solução de maior confiabilidade, eficiência computacional e previsibilidade geométrica para a resolução de sistemas lineares simétricos de alta complexidade.

\end{document}